% BHU PYQ -- Manually transcribed & error-corrected
% Paper: BCH-214 : Cost Accounting
% Exam: B.Com. (Hons.) Semester III Examination, 2022-23

\documentclass[11pt,a4paper]{article}
\usepackage[a4paper,top=2.5cm,bottom=2.5cm,left=2.8cm,right=2.8cm,headheight=14pt]{geometry}
\usepackage{amsmath,amssymb}
\usepackage[T1]{fontenc}
\usepackage[utf8]{inputenc}
\usepackage{lmodern,microtype,fancyhdr,enumitem,hyperref,lastpage,parskip,booktabs}

\hypersetup{
    colorlinks=true,
    urlcolor=black,
    linkcolor=black,
    pdftitle={BCH-214: Cost Accounting -- BHU B.Com. (Hons.) Sem III 2022-23}
}

\pagestyle{fancy}
\fancyhf{}
\renewcommand{\headrulewidth}{0.4pt}
\renewcommand{\footrulewidth}{0.4pt}
\fancyhead[L]{\small   \textit{Banaras Hindu University}}
\fancyhead[R]{\small   \textit{BCH-214: Cost Accounting}}
\fancyfoot[C]{\small Page  \thepage\ of \pageref{LastPage}}


\newlist{parts}{enumerate}{2}
\setlist[parts,1]{label=(\alph*),leftmargin=*,itemsep=5pt,topsep=3pt}
\setlist[parts,2]{label=(\roman*),leftmargin=*,itemsep=3pt,topsep=2pt}

\newcommand{\pts}[1]{\hfill   \textit{[#1]}}


\begin{document}


\begin{center}
  {\large\bfseries BANARAS HINDU UNIVERSITY}\\[2pt]
  {
\normalsize B.Com. (Hons.) --- Semester III Examination, 2022-23}\\[6pt]
  \rule{0.8\linewidth}{0.5pt}\\[6pt]
  {\large\bfseries Paper No. BCH-214: Cost Accounting}\\[4pt]
  {
\normalsize Time: 3 Hours\hfill Maximum Marks: 70}
\end{center}

\rule{\linewidth}{0.4pt}

\smallskip



\noindent   \textit{Instructions: Answer all questions. All questions carry equal marks (14 marks each).}

\bigskip




\noindent   \textbf{Question 1.} \\
It is said ``Cost Accounting is a system of foresight and not a postmortem examination, it turns losses into profits, speeds up activities and eliminates wastes''. Discuss the statement in detail.\pts{14}

\centerline{   \textbf{OR}}

How does Cost Accounting differ from Financial Accounting? Discuss the need of Cost Accounting in the present competitive industrial age.\pts{14}

\medskip\rule{\linewidth}{0.2pt}




\noindent   \textbf{Question 2.} \\
From the following data, prepare Stores Ledger on the basis of FIFO and LIFO Method:

\begin{itemize}
  \item   \textbf{January 03, 2023}: Received 5,000 units @ Rs. 9 p.u.
  \item   \textbf{January 04, 2023}: Issued 3,000 units
  \item   \textbf{January 06, 2023}: Issued 2,000 units
  \item   \textbf{January 07, 2023}: The weekly stock taking showed a shortage of 150 units.
  \item   \textbf{January 12, 2023}: Received 3,000 units @ Rs. 10 p.u.
  \item   \textbf{January 15, 2023}: Issued 1,500 units
\end{itemize}\pts{14}

\centerline{   \textbf{OR}}

A factory operates 50 weeks a year. The Diwali Bonus is fixed and is equal to 10 weeks' pay. The number of working hours in a week is 48 hours. The foreman receives a salary of Rs. 75 per week. There are 50 operators in the machine shop and 200 employees in the whole factory. The total power rating of the electric motors driving machines in the machine shop is 100 HP.

Compute Machine Hour Rate for a milling machine from the following details:

\begin{parts}
  \item The operator of the machine is paid Rs. 30 per week.
  \item The helper attending the machine is paid Rs. 15 per week.
  \item Repair expenses are Rs. 700.
  \item The purchase price of a milling machine is Rs. 13,000. The working life of the machine is estimated to be 24,000 hours. The estimated scrap value is Rs. 1,000.
  \item The power rating of the motor driving machine is 5 HP.
  \item The power consumption of the machine shop is Rs. 600 per month.
  \item Apart from the foreman's salary, machine shop overheads are Rs. 31,000 per year.
  \item The administrative overheads of the entire factory are Rs. 40,000 per year.
\end{parts}\pts{14}

\medskip\rule{\linewidth}{0.2pt}




\noindent   \textbf{Question 3.}

\begin{parts}
  \item A factory has received an order for three different types of casting weighing 18, 45 and 27 tons respectively. 10\% of the raw materials used are wasted in manufacturing and are sold as scrap for 20\% of the cost price of raw materials.
  
  The cost of raw materials is Rs. 300 per ton. The wages for three types of castings are Rs. 4,500, Rs. 11,000 and Rs. 5,500 respectively. The cost of the mould for the three different types of castings are Rs. 500, Rs. 600 and Rs. 400 respectively.
  
  If the factory overhead charges are 40\% of the prime cost in each case, find the cost of production per ton of each type of casting.\pts{10}
  
  \item What is material consumed? How is it calculated?\pts{4}
\end{parts}

\centerline{   \textbf{OR}}

The Manager of a cooler manufacturing concern consults you as a Cost Accountant as to the minimum price at which he can sell the cooler which are intended for mass production in future. The concern's records show the following particulars for the past year for the production and sale of 500 coolers:

\begin{itemize}
  \item   \textbf{Materials}: Rs. 1,20,000
  \item   \textbf{Direct wages}: Rs. 60,000
  \item   \textbf{Direct charges}: Rs. 10,000
  \item   \textbf{Works overhead}: Rs. 70,000
  \item   \textbf{Office overhead}: Rs. 28,000
  \item   \textbf{Selling overhead}: Rs. 32,000
  \item   \textbf{Profit}: Rs. 48,000
\end{itemize}

It is ascertained that 50\% of the work overhead fluctuate directly with production and 60\% of the selling overhead fluctuate with sales. It is anticipated that the concern would produce 2,500 coolers per annum and direct labour charge per unit will be reduced by 20\% while fixed overhead charge will increase by 25\%.

Prepare a Statement for Submission to the Manager of the concern if the same \% of profit is desired.\pts{14}

\medskip\rule{\linewidth}{0.2pt}




\noindent   \textbf{Question 4.} \\
A product passes through three distinct processes to completion. 10,000 units were introduced (Rs. 50,000) in Process A and the following expenses were incurred:


\begin{center}
  
\begin{tabular}{lrrr}
     oprule
   \textbf{Particulars} &   \textbf{Process A} &   \textbf{Process B} &   \textbf{Process C} \\
    \midrule
    Materials (Rs.) & 5,000 & 10,400 & 3,000 \\
    Wages (Rs.) & 25,000 & 35,000 & 26,000 \\
    Machine expenses (Rs.) & 6,000 & 5,000 & 3,600 \\
    Normal wastage & 5\% & 10\% & 5\% \\
    Scrap value per unit (Rs.) & 1.00 & 2.00 & 2.50 \\
    Output (units) & 9,000 & 8,000 & 7,700 \\
    ottomrule
  \end{tabular}
\end{center}

Prepare Process Account showing cost of output at each stage of manufacturing along with Normal Wastage Account, Abnormal Wastage Account and Abnormal Effectiveness Account.\pts{14}

\centerline{   \textbf{OR}}

Differentiate between by-products and joint products and explain the treatment of by-product in costing.\pts{14}

\medskip\rule{\linewidth}{0.2pt}




\noindent   \textbf{Question 5.} \\
A contractor took a contract for Rs. 15,00,000 on 1st January, 2021 and incurred the following expenses on the contract up to 30th June, 2022:

\begin{center}
  
\begin{tabular}{lr}
     oprule
   \textbf{Particulars} &   \textbf{Amount (Rs.)} \\
    \midrule
    Materials & 2,50,000 \\
    Wages & 3,00,000 \\
    Plant & 1,25,000 \\
    Overheads & 25,000 \\
    Wages outstanding & 37,500 \\
    Direct expenses & 25,000 \\
    Direct expenses outstanding & 12,500 \\
    Materials returned to stores & 5,000 \\
    Plant returned to stores (30th June, 2022) & 25,000 \\
    Plant stolen & 25,000 \\
    Materials lost & 12,500 \\
    Plant (costing Rs. 20,000) sold & 18,750 \\
    Cash received from contractee ($\frac{4}{5}$ of work certified) & 6,25,000 \\
    Cost of work uncertified & 50,000 \\
    Materials on hand & 12,500 \\
    ottomrule
  \end{tabular}
\end{center}

Depreciation on plant is 10\% per annum. Prepare Contract Account.\pts{14}

\centerline{   \textbf{OR}}

What are the essential features of Contract Costing? Explain the rules for the calculation of profit on incomplete contract.\pts{14}

\vfill
\rule{\linewidth}{0.4pt}

\begin{center}\small
  \href{https://bhu-exam.vercel.app/}{Click here to visit our website}\\[4pt]
  \href{https://me-aryan.vercel.app/}{Click here to visit our contributor}\\[4pt]
  \href{https://quashan-me.vercel.app/}{Click here to visit our contributor}
\end{center}

\end{document}
